\chapter{Methods}
\shorthandoff{-} 

\section{Data}

\subsection{Transcription factors}

The descriptive analysis is performed on three human transcription factors selected to span distinct binding modes relevant to the questions posed in this thesis: CTCF, MYC, and SPI1.
CTCF is an eleven-zinc-finger architectural factor with a long, well-characterised consensus motif and broad genomic occupancy, included as a strong-motif baseline against which embedding-based representations can be evaluated \cite{ongCTCFArchitecturalProtein2014}
MYC is a basic helix-loop-helix leucine-zipper factor that, in the analysis of \citet{faltejskovaNonconsensusFlankingSequence2026}, displays a pronounced directional dinucleotide asymmetry in the kilobase-scale flanks of its in vivo binding sites, making it a natural target for testing whether long-range sequence structure is reflected in Evo~2 representations \cite{luscherRegulationGeneTranscription2012a}.
SPI1 (PU.1) is an ETS-family lineage-specific factor that does not exhibit the directional flanking signature reported for MYC and is therefore included as a contrast factor for which the kilobase-scale signal is expected to be weaker or absent \cite{rothenbergMechanismsActionHematopoietic2019}.

\subsection{ChIP-seq data sources}

For each transcription factor, three ChIP-seq experiments are obtained from the ENCODE portal \cite{dunhamIntegratedEncyclopediaDNA2012}, yielding nine BED files in total \todo{cite ENCODE portal / ENCODE3 if a more recent reference is preferred}.
Conservative IDR-thresholded narrowPeak files \cite{liMeasuringReproducibilityHighthroughput2011} are used where available; otherwise, the IDR-thresholded peak file is used.
Cell-line selection follows the inclusion criterion adopted by \citet{faltejskovaNonconsensusFlankingSequence2026}, which retains only cell lines appearing in at least three experiments across their panel.
The accessions, transcription factors, and biosamples comprising the embedded set are listed in Table~\ref{tab:bedfiles}, alongside the number of peaks supplied to the embedding pipeline.
\begin{table}[h!]
  \begin{center}
    \caption{Chosen ChIP-seq experiments}
    \label{tab:table1}
    \begin{tabular}{l|c|r} % <-- Alignments: 1st column left, 2nd middle and 3rd right, with vertical lines in between
      \textbf{TF} & \textbf{Cell line} & \textbf{ENCODE accession}\\
      \hline
      MYC & GM12878 & ENCFF765CKK\\
      MYC & H1 & ENCFF700CXD\\
      MYC & MCF-7 & ENCFF149BIS\\
      CTCF & GM12878 & ENCFF017XLW\\
      CTCF & H1 & ENCFF692RPA\\
      CTCF & K562 & ENCFF769AUF\\
      SPI1 & GM12878 & ENCFF881ARS\\
      SPI1 & GM12878 & ENCFF198SPL\\
      SPI1 & K562 & ENCFF185DGL\\
    \end{tabular}
  \end{center}
\end{table}
Sequence is extracted from the human reference assembly hg19 to maintain consistency with \citet{faltejskovaNonconsensusFlankingSequence2026} and to permit direct comparison of per-bin readouts in the sanity check described in Section~\ref{sec:sanity}; this thesis does not consume any of the hg38-based annotation tracks (non-B DB, CpG islands, ENSEMBL regulatory features) used in the published study \todo{cite hg19 / GRCh37 (UCSC or Genome Reference Consortium)}.
BED files are downloaded programmatically from the ENCODE portal using the URL schema \texttt{https://www.encodeproject.org/files/\{accession\}/@@download/\{accession\}.bed.gz}.
\todo{fix wrapping?}

\subsection{Cell-line coverage}

The intended design covers the cell lines GM12878, H1, and K562 across all three transcription factors, providing a balanced grid for within-TF and across-TF comparisons.
The realised design departs from this target in one respect: the MYC K562 file (ENCFF043WTJ) was not embedded within the available compute window, and the MYC MCF-7 file (ENCFF149BIS) was substituted as the third MYC experiment.
The resulting cell-line coverage is therefore CTCF on GM12878, H1, K562; MYC on GM12878, H1, MCF-7; and SPI1 on GM12878 (two replicates) and K562.
The interpretive consequences of this imbalance for any TF--cell-line interaction analysis are addressed in the Results and Discussion chapters.

\subsection{Random-windows control}
\label{sec:randomctrl}

A random-windows control is constructed to provide a non-peak-centred reference that shares the window geometry, sequence preprocessing, and embedding pipeline of the ChIP-seq data and against which TF-centred positional structure can be contrasted.
Two thousand 10\,kb windows are sampled from hg19 autosomes and chromosome~X, with chromosome selection weighted by usable sequence length and chromosomes Y and M excluded.
Candidate windows are subjected to rejection sampling against the union of MYC and CTCF peak sets and against previously accepted control windows, ensuring the control is disjoint from those peak sets and internally non-overlapping.
SPI1 peaks are not included in the rejection set: the control was generated before SPI1 was added to the design and was not subsequently regenerated; given the small genomic footprint of the SPI1 peak set the residual chance of overlap is bounded but non-zero.
A fixed random seed of~42 is used throughout.

\subsection{Reproducibility and seeding}

Stochastic elements of the pipeline are seeded explicitly to support reproducibility of intermediate quantities.
A fixed random seed of~42 is applied at every sampling and stochastic-fitting site, including peak subsampling, random-windows generation, and UMAP fitting; this differs from the procedure of \citet{faltejskovaNonconsensusFlankingSequence2026}, in which peak subsampling is performed without a fixed seed.
The full source code, configuration files, and environment-setup scripts used to produce the embedded set and the analyses reported in subsequent chapters are available at the project repository \url{https://github.com/TypeTwo30312/tv-thesis}. 
\todo{probably footnote the repo} % \footnote{\url{...}}??

\section{Pipeline}

This part describes the deterministic chain from ENCODE BED files to per-peak embedding tensors stored on disk.
The three stages---sequence extraction, embedding extraction, and the computational environment in which they run---are each implemented as a separate Jupyter notebook in the project repository, and are reported here in the order in which they execute.

\subsection{Sequence extraction}
\label{sec:windows}

Window extraction reads an ENCODE narrowPeak file and outputs a per-peak record containing the peak coordinates and a fixed-length DNA sequence centred on the binding site.
The 10\,kb window length is chosen to match the $\pm 5000$\,bp range over which \citet{faltejskovaNonconsensusFlankingSequence2026} report statistically significant compositional structure around in vivo binding sites.
The peak midpoint is used as the binding-site proxy: for the ENCODE narrowPeak files inspected in this thesis, the reported summit coincides with the midpoint of the peak interval (CTCF GM12878 peaks are uniformly $\sim$356\,bp wide, MYC GM12878 peaks $\sim$390\,bp, MYC H1 $\sim$264\,bp), and the midpoint convention matches that used by \citet{faltejskovaNonconsensusFlankingSequence2026}.
Peaks whose 10\,kb window would extend past a chromosome boundary are dropped rather than padded, since padding with non-biological characters would introduce a fixed-position bias into any subsequent positional analysis.
Non-ACGT characters in the extracted sequence are masked to N at the extraction stage; the analysis-stage filtering of windows containing fully-N 50\,bp bins is described in Part~3 (Section~\ref{sec:nfilter}). \todo{adjust \ref{sec:nfilter} when Part~3 is drafted}
For each BED file, peaks are uniformly subsampled to a maximum of $10\,000$ (all peaks are retained if fewer are available), matching the \texttt{MAX\_SIZE\_SAMPLED} parameter used by \citet{faltejskovaNonconsensusFlankingSequence2026} but with a fixed random seed of~42 throughout.
Each output row carries the columns \texttt{peak\_id}, \texttt{chr}, \texttt{start}, \texttt{end}, \texttt{center}, and \texttt{sequence}, and is written as a gzip-compressed CSV.
The \texttt{ACCESSION\_\_TF\_\_BIOSAMPLE} naming convention used throughout the pipeline is preserved from the published source code released alongside \citet{faltejskovaNonconsensusFlankingSequence2026}, so that the per-bin feature-extraction scripts from that work remain directly applicable to the sanity check reported in Part~3 (Section~\ref{sec:sanity}). \todo{adjust \ref{sec:sanity} when Part~3 is drafted}

\subsection{Embedding extraction}
\label{sec:embeddings}

Embedding extraction is performed with the Evo~2 7B checkpoint released by \citet{brixiGenomeModellingDesign2026}; the architecture and pretraining objective are described in the background chapter (Section~\ref{sec:evo2bg}). \todo{adjust \ref{sec:evo2bg} to the actual background-chapter label}
The 7B family is used because it runs natively in bfloat16 on Ampere-class hardware and does not require Transformer Engine or FP8 support, which are needed for the 20B and 40B variants and are unavailable on the GPUs accessible through the CERIT-SC cluster used in this work. \todo{not true!! fix, reasoning}
Within the 7B family, the 1\,M-context variant (\texttt{evo2\_7b}) is selected over the 8\,kb-context base variant (\texttt{evo2\_7b\_base}) because the chosen 10\,kb window exceeds the latter's context length.
A single intermediate-layer embedding is extracted per token from the layer \texttt{blocks.28.mlp.l3}; intermediate layers are recommended over the final layer for representation learning because the final layer is specialised to the next-token prediction objective, and \texttt{blocks.28.mlp.l3} is the layer used by the exon-classifier example in the official Evo~2 repository \cite{brixiGenomeModellingDesign2026}. \todo{either run a comparison or suggest it in discussion}
Each 10\,kb sequence is tokenised under the Evo~2 byte-level scheme into a $(1, 10000)$ integer tensor (one token per nucleotide), passed through the model in a single forward pass on one GPU, and the resulting $(1, 10000, 4096)$ bfloat16 activation tensor is reshaped into 200 non-overlapping 50\,bp bins and averaged within each bin to yield a $(200, 4096)$ representation per peak, which is cast to fp16 before writing.
The 50\,bp bin width matches the $k$-mer resolution adopted by \citet{faltejskovaNonconsensusFlankingSequence2026} so that per-bin readouts are directly comparable; non-overlapping bins are used in place of their 25\,bp stride because each bin already integrates 50 input positions of the same forward pass and stride overlap would not contribute additional information.
Unlike the strictly local hand-engineered features used by \citet{faltejskovaNonconsensusFlankingSequence2026}, each binned embedding has the full 10\,kb window in its causal receptive field via the autoregressive forward pass, and so reflects model-internal aggregation of context up to that point rather than a fixed-window feature of the sequence itself.
A streaming-write strategy based on \texttt{numpy.memmap} is used for the per-file accumulator together with a checkpoint every 500 peaks, because retaining all per-peak tensors in memory before writing exceeds the per-pod RAM allocation on the cluster for the larger BED files.
The output of this stage is, per BED file, a single \texttt{.npz} archive containing the $(n_\text{peaks}, 200, 4096)$ fp16 array together with a \texttt{.parquet} sidecar carrying the \texttt{peak\_id} column and BED-derived metadata in the same row order.
Embedding extraction was performed on a single NVIDIA A40 GPU (46\,GB VRAM) with peak allocation of 19.17\,GB at batch size~1, at an observed throughput of approximately 0.5 peaks per second; the embedded set comprises the 9 ChIP-seq files of Part~1 plus the random-windows control of Section~\ref{sec:randomctrl}, totalling 96.9\,GB of compressed embeddings on disk.

\subsection{Computational environment}
\label{sec:env}

The embedding stage runs on the CERIT-SC JupyterHub platform under the NVIDIA PyTorch 2.5.0 container image, which provides a Python 3.10, CUDA 12.6, and Ubuntu 22.04 base; a project-specific Python 3.12 virtual environment is layered on top, installed via \texttt{uv} into the persistent home directory and pinning \texttt{torch}~2.7.1+cu128, \texttt{flash-attn}~2.8.0.post2, and \texttt{evo2}.
A separate CPU-only environment is used for the analyses reported in Part~3, isolating the heavy GPU stack from sessions that do not require it and pinning \texttt{numpy}, \texttt{polars}, \texttt{scikit-learn}, and \texttt{umap-learn} alongside the standard scientific Python packages.
Both environments are reproduced by setup scripts \todo{nope! single evo2 environment, no setup scripts} maintained alongside the source code, and the full code base, configuration files, and notebooks are available at the project repository\footnote{\url{https://github.com/TypeTwo30312/tv-thesis}}.

\section{Analyses}


\todo{probably the LDA, possibly the reverse sequence pass}
\shorthandon{-} 
%%%%%%%%%%%%%%%%%%%%%%%%%%%%%%%%%%%%

